\section*{Capítulo \ref{ch:g12:satellites-kepler} --- Satélites e movimento dos planetas}

\begin{solution}{exo:g12:satellites-kepler:1}
Na plataforma: $F = GMm/R^2 = \num{3.98e17}/\num{4.06e13} \approx
\qty{9.8e3}{N}$. A $r = \qty{6.77e6}{m}$: $F \approx \qty{8.7e3}{N}$
--- ainda $89\%$ do valor no solo. Os ocupantes flutuam porque
\emph{caem} junto com a estação, e não porque a gravidade sumiu.
\end{solution}

\begin{solution}{exo:g12:satellites-kepler:2}
$v = \sqrt{GM/R} = \sqrt{\num{3.98e14}/\num{6.37e6}} \approx
\qty{7.9e3}{m/s}$; $T = 2\pi R/v \approx \qty{5.1e3}{s} \approx
\qty{84}{min}$. Exatamente os \qty{7.9}{km/s} da bala de canhão: uma
\omterm{def:g10:universal-gravitation:satellite}{órbita} rasante \emph{é} o tiro de Newton.
\end{solution}

\begin{solution}{exo:g12:satellites-kepler:3}
$v \propto 1/\sqrt{r}$: dividida por $\sqrt 2 \approx 1.41$.
$T \propto r^{3/2}$: multiplicado por $2^{3/2} \approx 2.8$. Nenhum dos
dois contém $m$: não.
\end{solution}

\begin{solution}{exo:g12:satellites-kepler:4}
$v = \sqrt{\num{3.98e14}/\num{6.77e6}} \approx \qty{7.67e3}{m/s}$;
$T = 2\pi r/v \approx \qty{5.55e3}{s} \approx \qty{92}{min}$;
$\num{86400}/5546 \approx 15.6$ \omterm{def:g10:universal-gravitation:satellite}{órbitas} por dia.
\end{solution}

\begin{solution}{exo:g12:satellites-kepler:5}
\omterm{def:g12:satellites-kepler:ellipse}{Elipses} com o Sol em um dos \omterm{def:g12:satellites-kepler:ellipse}{focos}; áreas iguais em tempos iguais;
$T^2/a^3$ idêntico para todos os planetas. Mais rápido no \omterm{def:g12:satellites-kepler:ellipse}{periélio}, pela
segunda lei (o segmento curto que varre precisa varrer depressa). Não: a
massa do \omterm{def:g10:universal-gravitation:satellite}{satélite} se cancela --- só entra a massa central.
\end{solution}

\begin{solution}{exo:g12:satellites-kepler:6}
$T = \qty{2.36e6}{s}$: $M = 4\pi^2 r^3/(GT^2) = \num{2.24e27}/371
\approx \qty{6.0e24}{kg}$ --- a menos de $1\%$ de \qty{5.97e24}{kg}.
\end{solution}

\begin{solution}{exo:g12:satellites-kepler:7}
A gravidade aponta para o centro da Terra e, no movimento circular, a
\omterm{def:g11:forces-and-motion:netforce}{força resultante} aponta para o centro do círculo: os dois centros
coincidem. Um círculo acima do paralelo de Paris seria centrado no eixo,
e não no centro da Terra --- impossível. Pairar só funciona onde o
paralelo é um círculo máximo centrado no centro: sobre o \emph{equador}.
\end{solution}

\begin{solution}{exo:g12:satellites-kepler:8}
$T = \qty{2.754e4}{s}$: $M = 4\pi^2 r^3/(GT^2) =
\num{3.26e22}/\num{5.06e-2} \approx \qty{6.4e23}{kg}$ --- o Marte da
ficha de dados.
\end{solution}

\begin{solution}{exo:g12:satellites-kepler:9}
Terra: $T^2/a^3 = \num{9.96e14}/\num{3.35e33} = \qty{2.98e-19}{s^2/m^3}$.
Marte: $T = \qty{5.94e7}{s}$, $\num{3.52e15}/\num{1.18e34} =
\qty{2.98e-19}{s^2/m^3}$ --- iguais, como exige a terceira lei. Então
$M = 4\pi^2/(G \times \num{2.98e-19}) \approx \qty{1.99e30}{kg}$.
\end{solution}

\begin{solution}{exo:g12:satellites-kepler:10}
$v = 2\pi r/T = \qty{1.02e3}{m/s}$; $a = v^2/r = \qty{2.72e-3}{m/s^2}$;
$g/60^2 = 9.81/3600 = \qty{2.73e-3}{m/s^2}$. A Lua está a $60$ raios
terrestres, e a queda dela é $60^2$ vezes mais fraca do que a da maçã:
exatamente o inverso do quadrado --- uma lei só para as duas.
\end{solution}

\begin{solution}{exo:g12:satellites-kepler:11}
Ceres: $T = 2.77^{3/2} = \sqrt{21.3} \approx \qty{4.6}{yr}$. Halley:
$T = 17.8^{3/2} \approx \qty{75}{yr}$, de modo que o próximo \omterm{def:g12:satellites-kepler:ellipse}{periélio}
fica em torno de $1986 + 75 \approx 2061$.
\end{solution}

\begin{solution}{exo:g12:satellites-kepler:12}
A \qty{400}{km}: $v \approx \qty{7.67e3}{m/s}$; a \qty{350}{km}
($r = \qty{6.72e6}{m}$): $v \approx \qty{7.70e3}{m/s}$. O \omterm{def:g10:inertia:inventory}{atrito} de fato
drena \omterm{def:g11:mechanical-energy:em}{energia mecânica} --- a \omterm{def:g10:universal-gravitation:satellite}{órbita} encolhe; mas, na descida, o \omterm{def:g11:work-of-force:work}{trabalho}
da gravidade paga com juros a velocidade perdida para o arrasto: $E_m$
cai enquanto $E_k$ sobe.
\end{solution}

\begin{solution}{exo:g12:satellites-kepler:13}
$T = \qty{3.024e5}{s}$: $M = 4\pi^2 r^3/(GT^2) = \num{1.65e31}/6.10
\approx \qty{2.7e30}{kg} \approx 1.4$ Sóis. A massa do \emph{planeta}
está fora de \omterm{def:g12:projectile-motion:range}{alcance}: ela se cancelou nas equações da \omterm{def:g10:universal-gravitation:satellite}{órbita}.
\end{solution}

\begin{solution}{exo:g12:satellites-kepler:14}
$T = \qty{8.862e4}{s}$, $GM = \qty{4.28e13}{m^3/s^2}$:
$r^3 = GMT^2/(4\pi^2) = \num{8.52e21}$, logo $r \approx \qty{2.04e7}{m}$
e $h = r - R \approx \qty{1.7e7}{m} \approx \qty{17000}{km}$ acima do
solo marciano.
\end{solution}

\begin{solution}{exo:g12:satellites-kepler:15}
$v_a = r_p v_p / r_a = \num{8.77e10} \times \num{5.45e4} /
\num{5.25e12} \approx \qty{9.1e2}{m/s}$. Como $r_a/r_p \approx 60$, o
cometa está $60$ vezes mais longe e vai $60$ vezes mais devagar: ele
passa décadas se arrastando no escuro e semanas disparando ao redor do
Sol.
\end{solution}

\begin{solution}{pb:g12:satellites-kepler:1}
\textbf{1.} $T = \qty{86164}{s}$: o solo gira com a Terra \emph{em
relação às estrelas} uma vez por \omterm{def:g12:satellites-kepler:geo}{dia sideral}; o dia de \qty{24}{h}
acrescenta o grau por dia que a Terra avança em torno do Sol.

\textbf{2.} A \omterm{def:g11:forces-and-motion:netforce}{força resultante} do movimento circular mira o centro do
círculo; a única \omterm{def:g10:inertia:force}{força}, a gravidade, mira o centro da Terra: os dois
centros são o mesmo ponto.

\textbf{3.} Um círculo pairando sobre Paris seria centrado no eixo, ao
norte do centro --- excluído por 2. A \omterm{def:g10:universal-gravitation:satellite}{órbita} precisa estar no
\emph{plano equatorial}.

\textbf{4.} Componente normal da \omterm{thm:g12:newtons-laws:second}{segunda lei de Newton}:
$m v^2/r = GmM/r^2$, logo $v = \sqrt{GM/r}$ ($m$ se cancela).

\textbf{5.} $T = 2\pi r/v = 2\pi\sqrt{r^3/(GM)}$; elevando ao quadrado,
$T^2/r^3 = 4\pi^2/(GM)$.

\textbf{6.} $r = \left(GMT^2/4\pi^2\right)^{1/3} = (\num{7.49e22})^{1/3}
\approx \qty{4.22e7}{m}$.

\textbf{7.} $h = r - R = \num{4.22e7} - \num{6.37e6} \approx
\qty{3.58e7}{m} \approx \qty{35800}{km}$.

\textbf{8.} $v = 2\pi \times \num{4.22e7}/\num{86164} \approx
\qty{3.07e3}{m/s}$.

\textbf{9.} $\sqrt{\num{3.98e14}/\num{4.22e7}} \approx
\qty{3.07e3}{m/s}$ --- coerente.

\textbf{10.} Raios idênticos: a massa se cancelou em 4. O \omterm{def:g10:universal-gravitation:satellite}{satélite} mais
pesado muda apenas a conta do lançamento --- mais combustível para a
mesma \omterm{def:g10:universal-gravitation:satellite}{órbita}.

\textbf{11.} $v = \sqrt{\num{3.98e14}/\num{6.77e6}} \approx
\qty{7.67e3}{m/s}$.

\textbf{12.} $T = 2\pi r/v \approx \qty{5.55e3}{s} \approx
\qty{92.4}{min}$.

\textbf{13.} $\num{86400}/5546 \approx 15.6$ \omterm{def:g10:universal-gravitation:satellite}{órbitas}: cerca de $16$
nasceres do Sol por dia.

\textbf{14.} $r_{\text{geo}}/r_{\text{EEI}} = \num{4.22e7}/\num{6.77e6}
= 6.23$; $6.23^{3/2} \approx 15.5 = \num{86164}/5546$ --- a terceira lei
de Kepler, verificada na hora.

\textbf{15.} A estação é mais rápida (\qty{7.67}{km/s} contra
\qty{3.07}{km/s}), e ainda assim o \omterm{def:g12:satellites-kepler:geo}{geoestacionário} custou mais \omterm{def:g10:energy-conservation:energy}{energia}
por \omterm{def:g10:orders-of-magnitude:unit}{quilograma}: a subida em \omterm{def:g10:energy-conservation:energy}{energia} potencial supera a perda de
velocidade.

\textbf{16.} $M_J = 4\pi^2 r^3/(G T^2)$.

\textbf{17.} $T = \qty{1.529e5}{s}$: $M_J = \num{2.97e27}/1.56 \approx
\qty{1.90e27}{kg}$.

\textbf{18.} $\num{1.90e27}/\num{5.97e24} \approx 318$ Terras --- e
cerca de $1/1000$ do Sol.

\textbf{19.} $T_J = \qty{3.576e4}{s}$: $r^3 = GM_J T_J^2/(4\pi^2) =
\num{4.11e24}$, $r \approx \qty{1.60e8}{m}$ --- Io, a \qty{4.22e8}{m},
voa $2.6$ vezes mais alto e deriva pelo céu de Júpiter.

\textbf{20.} Estacione o repetidor no círculo equatorial de altitude
\qty{35800}{km}, onde ele corre a \qty{3.07}{km/s} e nunca sai da mira
da antena; e a mesma lei, alimentada com o mês de Io, pesou Júpiter em
$\qty{1.90e27}{kg}$ --- $318$ Terras --- de graça.
\end{solution}
